\documentclass{article}
\usepackage{polyglossia}
\setmainlanguage{english}
\setotherlanguage{japanese}
\setotherlanguage{chinese}
\setmainfont{Latin Modern Roman}
\newfontfamily\japanesefont{Latin Modern Roman}
\newfontfamily\chinesefont{Latin Modern Roman}
\begin{document}
English text first.
\textjapanese{「こんにちは」、世界。日本語の文章です。句読点、括弧（かっこ）や「かぎ括弧」も。ちょっと待ってください……ねえ、ゃゅょっァィゥェォ。東京は２０２６年１０月６日（火）です。}
Back to English.

\begin{japanese}
吾輩は猫である。名前はまだ無い。どこで生れたかとんと見当がつかぬ。何でも薄暗いじめじめした所でニャーニャー泣いていた事だけは記憶している。
\end{japanese}

\textchinese{这是一个测试：中文、英文（English）和数字123混排。《三国演义》第一回——宴桃园豪杰三结义。}
\end{document}
